\documentclass[12pt,a4paper]{article}

\usepackage[utf8]{inputenc}
\usepackage[T1]{fontenc}
\usepackage[margin=2.5cm]{geometry}
\usepackage{hyperref}
\usepackage{microtype}
\usepackage{parskip}
\usepackage{lmodern}

\hypersetup{colorlinks=true, urlcolor=blue!60!black}

\pagestyle{empty}

\begin{document}

\begin{flushright}
Kundai Farai Sachikonye \\
Auf der Lichtung 19, 82178 Puchheim, Germany \\
\href{https://github.com/fullscreen-triangle}{github.com/fullscreen-triangle} \\
\texttt{kundai.sachikonye@bitspark.com} \\
(+49) 1626386344 \\[6pt]
18 May 2026
\end{flushright}

\vspace{1em}

Dr.\ Moritz Kleinert \\
Photonic Biosensors \\
Fraunhofer-Institut für Nachrichtentechnik --- Heinrich-Hertz-Institut (HHI) \\
Berlin \\
Ref.\ 84447

\vspace{1.5em}

\textbf{Re: Wissenschaftliche*r Mitarbeiter*in für Biochemie und Oberflächenfunktionalisierung (Ref.\ 84447)}

\vspace{1em}

Dear Dr.\ Kleinert,

I am applying for the scientist position in biochemistry and surface
functionalization at Fraunhofer HHI (Ref.\ 84447).
My background combines formal biochemistry training and practical
protein-handling experience with independent computational work on the
optical physics and molecular recognition problems that underlie photonic
biosensor performance.
The surface chemistry specific to photonic integrated chips is the aspect
I have not done on the bench; the biochemical foundation, the understanding
of evanescent-field sensing mechanisms, and the analytical tools for
biosensor signal processing are in place.

\textbf{Biochemistry and assay development background.}
My BSc in Biochemistry and Cell Biology (Jacobs University Bremen) and
MSc in Pharmaceutical Biotechnology (HAW Hamburg) together cover the
full biochemical toolkit the position requires: protein purification and
handling, antibody-based detection assays, bioconjugation strategies,
ELISA development, and quality and reproducibility standards in bioanalytical
method development.
At the Universitätsklinikum Hamburg-Eppendorf (2017--2018) I developed
automated LC-MS/MS and MALDI-TOF pipelines for proteomics and glycomics ---
work that required systematic protein immobilisation on MALDI targets,
surface preparation, and strict protocol standardisation to achieve
reproducible ion signals across sample series.
At the University Medical Center Groningen (2013) I performed UPLC-MS/MS
shotgun proteomics for cancer biomarker discovery, which required developing
sample preparation and digestion protocols validated against a reference
standard. These are the skills of assay development and experimental
validation; the substrate here is PIC chip surfaces rather than MS plates,
but the biochemical logic --- controlled surface chemistry, selective
capture, non-specific binding suppression, reproducibility at the device
level --- is the same.

\textbf{Optical physics and photonic biosensing mechanisms.}
The detection mechanism in label-free photonic biosensing is a refractive
index change in the evanescent field at the waveguide surface.
Understanding where that signal comes from --- and how large it should be
for a given analyte-receptor interaction --- requires first-principles optics.
The \href{https://github.com/fullscreen-triangle/lagrangian}{Lagrangian}
platform I developed proves Beer-Lambert absorption, Kirchhoff circuit laws,
and thermal emission as formally identical operations via the Resonant Stack
theorem: transmittance $T = 1/e = 0.3679$ is the universal fixed point in all
three representations. This is exactly the physics of evanescent field coupling
and absorption in a waveguide: the same partition-geometry argument that governs
bulk absorption governs the sensing layer interaction.
The \href{https://honjo-masamune.vercel.app/tools}{Borgia} framework derives
electronic transitions, absorption, and emission spectra of molecules from
bounded phase space geometry alone, with zero empirical parameters.
Validated against NIST for nine elements and six molecules (self-consistency
errors $<$1.63\,\%). For photonic biosensor design, this means I can predict
the absorption cross-section and refractive index contribution of surface-bound
molecules from molecular structure, before running the experiment --- useful for
pre-screening linker chemistries and bioreceptor candidates.

\textbf{Molecular recognition and receptor--analyte binding.}
The quantitative output of a label-free biosensor is the binding kinetics
and equilibrium $K_d$ of the receptor--analyte interaction.
The \href{https://github.com/fullscreen-triangle/crown-prince}{Crown Prince}
framework derives dissociation constants from bounded phase space geometry:
$K_d$ is determined by the partition structure of the receptor--ligand
bounded oscillatory system, without empirical fitting.
Validated against 39 NIST compounds.
The \href{https://syndrome-seven.vercel.app/tools/mhc-binding}{Syndrome}
platform implements spectral DFT-based molecular recognition: three
physicochemical channels (hydropathy, volume, charge) → Fourier transform
→ 36-dimensional cosine similarity $\rho$; binding is predicted when
$\rho \geq \rho^* = 0.72$ with $O(cL\log L)$ complexity.
The mathematical structure is identical to what a photonic biosensor measures:
spectral complementarity between receptor and analyte, with a threshold
separating bound from unbound states.
This gives me a computational layer for receptor design and binding affinity
estimation that complements experimental validation on the chip.

\textbf{Signal analysis and data evaluation.}
The \href{https://github.com/fullscreen-triangle/lavoisier}{Lavoisier}
framework processes spectral data at $>$100\,GB scale: dual-pipeline numerical
and CNN-based analysis (PyTorch), 98.9\,\% NIST feature-extraction accuracy,
Ray distributed computing. The signal analysis problem in photonic biosensing
(extracting binding signal from interferometric or resonance-shift data against
background noise and matrix interference) is structurally identical to the
spectral peak extraction and ion annotation problem in mass spectrometry.
I can therefore bring both analytical rigour and practical tooling to the
data evaluation component of the role.

\textbf{Background.}
I hold a BSc in Biochemistry and Cell Biology (Jacobs University Bremen),
an MSc in Pharmaceutical Biotechnology (HAW Hamburg), and an MSc in
Computational Biology (Universität Tübingen), and I conducted doctoral
research in Computational Lipidomics at TUM.
I am legally resident in Germany with full work authorisation and available
immediately, including relocation to Berlin.
English is native; German is conversational.

\vspace{1.5em}
Yours sincerely,

\vspace{2.5em}
\textbf{Kundai Farai Sachikonye}

\end{document}
